\documentclass{article}
\usepackage[utf8]{inputenc}

\title{Relatório: História do Python}
\author{João Francisco Silva Correia}
\date{13/08/2026}

\begin{document}

\maketitle

\section{Introdução}.
O Python foi criado por Guido van Rossum como uma linguagem interpretada e focada na legibilidade do código[cite: 1].

\begin{itemize}
    \item \textbf{Git}: Utilizado para o controle de versão[cite: 1].
    \item \textbf{Make}: Ferramenta para automação de tarefas de compilação[cite: 1].
    \item \textbf{Sed}: Utilitário para filtragem e transformação de textos via terminal[cite: 1].
\end{itemize}.

\section{Conclusão}
Esta é a conclusão do relatório falando sobre a historia do pyhton.

A filosofia do Python enfatiza simplicidade e legibilidade no desenvolvimento de software \cite{python}[cite: 1].

\bibliographystyle{plain}
\bibliography{referencias}

\end{document}
